\section{Introduction}
\label{sec:introduction}

Thermal energy storage (TES) has emerged as a critical technology for addressing the growing mismatch between energy supply and demand in applications ranging from renewable energy harvesting to electronics thermal management~\cite{mehrpooya2026phase,tyagi2012phase,mckenna2021thermal}. Among the various TES approaches, latent heat thermal energy storage (LHTES) using phase change materials (PCMs) offers the distinct advantage of storing and releasing large quantities of thermal energy at nearly constant temperatures, corresponding to the material's phase transition point~\cite{sait2022solidification,belusko2018direct}. Paraffin waxes, in particular, have attracted considerable attention due to their high latent heat capacity, chemical stability, negligible supercooling, and relatively low cost~\cite{bharathiraja2024thermal,bharathiraja2023studies,lin2014thermal,abhat1983low}. The thermophysical properties and application domains of paraffin-based PCMs have been surveyed extensively~\cite{veerakumar2016phase}, and recent work has highlighted the impact of system size and thermal gradients on supercooling behavior~\cite{lilley2020impact}.

Despite these favorable properties, the practical application of paraffin-based PCMs is severely limited by their inherently low thermal conductivity (typically 0.15--0.35~W/(m$\cdot$K)), which results in sluggish charging and discharging rates~\cite{sciacovelli2012melting,amin2014effective}. To overcome this limitation, numerous thermal conductivity enhancement techniques have been proposed, including the incorporation of high-conductivity fillers such as metal foams~\cite{ghalambaz2020conjugate,parida2020effect,liu2024pore}, extended surfaces (fins)~\cite{patel2018thermal,temel2023experimental,altinok2026asymmetric,uniyal2024melting}, nanoparticles~\cite{murugan2017thermal,bashirpour2021investigation,bora2026enhanced}, and encapsulation strategies~\cite{solomon2016exergy,alkayiem2012thermal,sulzgruber2021micro}. In parallel, form-stable composite PCMs, in which paraffin is confined within expanded graphite or polymer matrices, simultaneously improve thermal conductivity and prevent leakage~\cite{sari2007thermal,li2013nano,xia2010preparation,zhang2006study,zhai2017preparation,zuo2020enhanced,guler2026engineered,zhang2025sulfur,sari2004form}. While each approach has demonstrated measurable improvements, the development of structurally integrated heat transfer enhancement remains an active frontier.

Triply periodic minimal surfaces (TPMS) represent a class of mathematically defined surfaces that partition three-dimensional space into two continuous, interpenetrating channels with zero mean curvature~\cite{jung2006fluid,himmelmann2026amorphous}. The most prominent TPMS families---Gyroid, Primitive (Schwarz-P), Diamond (Schwarz-D), and IWP---possess exceptional geometric properties including high surface-area-to-volume ratios, smooth curvature transitions, and tunable porosity~\cite{tang2023analysis,kaur2021flow,barakat2024controlling}. These characteristics have made TPMS structures increasingly attractive for heat transfer applications, with recent studies demonstrating their superior convective heat transfer performance compared to conventional geometries~\cite{tang2023experimental,renon2025numerical,chen2026comprehensive}. Critically, advances in metal additive manufacturing---particularly laser powder bed fusion (L-PBF)---have enabled the fabrication of complex TPMS lattice structures in high-conductivity metals such as aluminum alloys~\cite{catchpole2019thermal,chouhan2025heat,fouad2026additive,noronha2025thin}, opening the possibility of using TPMS lattices as both structural scaffolds and thermal conductivity enhancers for PCM composites.

The convergence of TPMS lattice design, additive manufacturing, and PCM-based thermal storage represents a nascent but highly promising research direction. Recent experimental work by Piacquadio et al.~\cite{piacquadio2022experimental} demonstrated that PCM embedded in additively manufactured lattice structures exhibits significantly enhanced thermal response compared to bare PCM. However, systematic investigations comparing different TPMS topologies (Gyroid vs.\ Primitive vs.\ IWP) under varying heating conditions remain scarce, and the thermal performance of such composites during both melting and solidification has not been comprehensively characterized.

\emph{First}, we present an experimental investigation of paraffin-based PCM (phase change temperature 42\textdegree C) infiltrated into AlSiMg TPMS lattice structures fabricated by L-PBF, with systematic characterization under Gyroid topology at three heating powers (10~W, 20~W, 30~W).

\emph{Second}, we establish the spatial temperature distribution characteristics within the TPMS/PCM composite using a nine-sensor array organized into three groups, revealing distinct thermal stratification patterns.

\emph{Third}, we lay the experimental foundation for a comparative study across three TPMS topologies (Gyroid, Primitive, and IWP), with the Gyroid results presented here serving as the baseline for forthcoming Primitive and IWP experiments.

The remainder of this paper is organized as follows. Section~\ref{sec:related} reviews related work on PCM thermal conductivity enhancement and TPMS heat transfer. Section~\ref{sec:method} describes the TPMS lattice design, PCM selection, and additive manufacturing process. Section~\ref{sec:experiment} details the experimental setup and measurement protocol. Section~\ref{sec:results} presents and analyzes the temperature evolution data. Section~\ref{sec:conclusion} summarizes the key findings and outlines future work.
